\documentclass[11pt]{article}

\usepackage[margin=1in]{geometry}
\usepackage{amsmath}
\usepackage{booktabs}
\usepackage{longtable}
\usepackage{tabularx}
\usepackage{array}
\usepackage{enumitem}
\usepackage{hyperref}
\usepackage{xcolor}
\usepackage{microtype}

\setlist[itemize]{leftmargin=*}
\setlist[enumerate]{leftmargin=*}
\setlength{\parskip}{0.5em}
\setlength{\parindent}{0pt}

\newcolumntype{Y}{>{\raggedright\arraybackslash}X}

\title{\textbf{DNA Paper Roadmap}\\Second-Pass Publication Plan}
\author{Recognition Science DNA Program}
\date{March 2026}

\begin{document}
\maketitle

\section*{Status}
Second-pass publication roadmap after the March 2026 Lean DNA buildout.

This document is the paper plan, not the open-problems ledger. Its purpose is to decide:
\begin{itemize}
\item which DNA papers are strongest now,
\item which claims are already supported by the Lean surface,
\item which papers need one more computational or empirical pass,
\item and what order gives the best mix of scientific impact, clarity, and publishability.
\end{itemize}

\section*{What Changed Since The First Pass}
The main constraint is no longer whether the theory can be formalized at all. A substantial Tier 2 / Tier 3 Lean layer now exists and \texttt{IndisputableMonolith.Genetics} builds cleanly.

The new backbone includes:
\begin{itemize}
\item reusable cost and certificate infrastructure:
\begin{itemize}
\item \texttt{MutationModel.lean}
\item \texttt{CompositeJCost.lean}
\item \texttt{OptimizationCertificate.lean}
\end{itemize}
\item strengthened Tier 2 theory packages:
\begin{itemize}
\item \texttt{IntronPlacementOptimality.lean}
\item \texttt{GearboxSelectivity.lean}
\item \texttt{ZParityDegeneracyTable.lean}
\item \texttt{WobbleZParity.lean}
\item \texttt{ConservationHierarchy.lean}
\end{itemize}
\item Tier 3 theorem shells and certificate interfaces:
\begin{itemize}
\item \texttt{CodonAssignmentObjective.lean}
\item \texttt{CodonAssignmentOptimality.lean}
\item \texttt{WTokenAssignmentBridge.lean}
\item \texttt{DegeneracySearchSpace.lean}
\item \texttt{CodonAssignmentGlobalOptimality.lean}
\item \texttt{SyntheticFalsifierAssays.lean}
\end{itemize}
\end{itemize}

This changes the paper program materially. The main bottleneck now is evidence packaging, deterministic computation, and reviewer-facing framing, not missing Lean scaffolding.

\section*{Executive Summary}

\subsection*{Tier A: Draft Now}
\begin{enumerate}
\item \textbf{The standard genetic code is not locally optimal under simple chemical distance}
\item \textbf{Z-parity and the hierarchy of conservation in the genetic code}
\item \textbf{Toward a zero-junk genome} (revise and reissue as the anchor paper)
\item \textbf{Sex as phase-debt cancellation} (draft once one clean toy model is added)
\end{enumerate}

\subsection*{Tier B: One More Computational Or Validation Pass}
\begin{enumerate}[start=5]
\item \textbf{Intron positioning from helical phase stress}
\item \textbf{The molecular gearbox: phi-ladder timing in biology}
\item \textbf{Synthetic biology falsifiers for the RS DNA framework} as a design / assay paper
\end{enumerate}

\subsection*{Tier C: Flagship Long-Horizon Papers}
\begin{enumerate}[start=8]
\item \textbf{Deriving the exact codon-to-amino-acid assignment}
\item \textbf{Global J-cost optimality of the genetic code}
\end{enumerate}

\section*{Portfolio At A Glance}
\footnotesize
\begin{longtable}{@{}p{0.03\textwidth}p{0.22\textwidth}p{0.12\textwidth}p{0.14\textwidth}p{0.22\textwidth}p{0.19\textwidth}@{}}
\toprule
\# & Paper & Readiness now & Lean status & Main blocker & Best role \\
\midrule
1 & Code not chemically optimal & Write now & Strong & Publication-grade benchmark package & Sharp computational result \\
2 & Sex as phase-debt cancellation & Near-write & Strong core & Standard genetics translation + toy simulation & Conceptual theory paper \\
3 & Zero-junk genome & Revise now & Strong & Figures and tighter framing & Anchor / platform paper \\
4 & Intron positioning & One more pass & Theory-ready & Predictor + gene benchmark & Genomics paper \\
5 & Molecular gearbox & One more pass & Theory-ready & Mechanism tightening + spectroscopy review & Biophysics bridge paper \\
6 & Z-parity hierarchy & Write now & Very strong & Interpretation + figures & Small clean paper \\
7 & Exact codon assignment & Not yet claim-ready & Interface-ready & Composite metric calibration + local certificates & Flagship derivation \\
8 & Global J-cost optimality & Long-horizon & Interface-ready & Search reduction + certified computation & Flagship theorem / computation \\
9 & Synthetic falsifiers & Design-paper ready & Interface-ready & Construct definitions + assay path + collaborators & Experimental strategy paper \\
\bottomrule
\end{longtable}
\normalsize

\section*{Shared Paper Rules}
Every DNA paper should obey the same discipline:
\begin{itemize}
\item \texttt{Theorem}: cite a named Lean result and the exact module that contains it.
\item \texttt{Model}: clearly label abstractions formalized as definitions / interfaces but not yet experimentally anchored.
\item \texttt{Hypothesis}: explicitly state what would falsify the claim.
\item \texttt{Computation}: ship deterministic scripts, frozen outputs, and result tables.
\item \texttt{Scope}: do not claim ``RS proves X biologically'' when Lean only proves the structural shell.
\end{itemize}

Every paper package should also include:
\begin{itemize}
\item one paragraph explaining what is formally proved,
\item one paragraph explaining what remains model-level,
\item one reviewer-facing comparison to the closest standard literature,
\item a methods appendix with scripts and dataset versions,
\item and a table mapping major claims to theorem / model / hypothesis.
\end{itemize}

\section*{Current Lean Baseline For The Paper Program}

\subsection*{Codon objective and optimization}
\begin{itemize}
\item \texttt{CodonNeighborCost}
\item \texttt{ErrorImpactMatrix}
\item \texttt{MutationModel}
\item \texttt{CompositeJCost}
\item \texttt{CodonAssignmentObjective}
\item \texttt{CodonAssignmentOptimality}
\item \texttt{CodonAssignmentGlobalOptimality}
\item \texttt{DegeneracySearchSpace}
\end{itemize}
This is the backbone for Papers 1, 8, and 9.

\subsection*{Wobble, parity, and conservation}
\begin{itemize}
\item \texttt{AminoAcidZClassification}
\item \texttt{ZInvariantPipeline}
\item \texttt{ZParityDegeneracyTable}
\item \texttt{WobbleZParity}
\item \texttt{ConservationHierarchy}
\end{itemize}
This is the backbone for Paper 6 and part of Paper 3.

\subsection*{Helical phase and intron placement}
\begin{itemize}
\item \texttt{IntronStructure}
\item \texttt{PhaseAlignmentLayer}
\item \texttt{HelicalPhaseModel}
\item \texttt{IntronPlacementOptimality}
\end{itemize}
This is the backbone for Paper 5 and part of Paper 3.

\subsection*{Timing, FOLD, and gearbox}
\begin{itemize}
\item \texttt{TemporalCoordinationLayer}
\item \texttt{ChromosomePhaseProfile}
\item \texttt{HydrationGearbox}
\item \texttt{GearboxSelectivity}
\end{itemize}
This is the backbone for Papers 2, 6, and part of Paper 3.

\subsection*{Bridge layer to WTokens and synthetic tests}
\begin{itemize}
\item \texttt{WTokenAssignmentBridge}
\item \texttt{SyntheticFalsifierAssays}
\item \texttt{OptimizationCertificate}
\end{itemize}
This is the backbone for Papers 7, 8, and 9.

All of the above modules are currently \texttt{sorry}-free and \texttt{axiom}-free.

\section*{Detailed Paper Plans}

\subsection*{1. The Standard Genetic Code Is Not Locally Optimal Under Simple Chemical Distance}
\textbf{Why this should be written now.} This is still the sharpest tactical result in the program. It is concrete, computational, reproducible, and easy to explain: under a naive chemistry-only metric, the standard code is not locally optimal, so the true biological objective must be richer.

\textbf{Lean base now in place.}
\begin{itemize}
\item \texttt{AminoAcidZClassification}
\item \texttt{CodonNeighborCost}
\item \texttt{ErrorImpactMatrix}
\item \texttt{MutationModel}
\item \texttt{CompositeJCost}
\item \texttt{CodonAssignmentObjective}
\item \texttt{CodonAssignmentOptimality}
\end{itemize}

\textbf{Strongest honest claim.} The standard genetic code does not minimize a simple chemical-distance objective. Any serious optimality theory of the code must include additional structure, and RS gives a principled place to look for it.

\textbf{Still needed before submission.}
\begin{itemize}
\item deterministic optimizer reruns
\item frozen result tables
\item random-code baseline
\item degeneracy-preserving shuffle baseline
\item pairwise-swap histogram
\item annealing convergence figure
\item supplement listing the best improving swaps
\end{itemize}

\textbf{Suggested title.} \emph{The Standard Genetic Code Is Not Locally Optimal Under Simple Chemical Distance}

\textbf{Priority.} Highest.

\subsection*{2. Sex as Phase-Debt Cancellation}
\textbf{Why this is still high value.} This remains the strongest conceptual biology result outside the zero-junk platform paper. It gives RS a structural answer to why recombination helps.

\textbf{Lean base now in place.}
\begin{itemize}
\item \texttt{TemporalCoordinationLayer}
\item \texttt{ChromosomePhaseProfile}
\end{itemize}

\textbf{Strongest honest claim.} The FOLD operator annihilates antisymmetric phase debt and asexual inheritance does not. This gives a formal structural analogue of recombination advantage.

\textbf{Still needed before submission.}
\begin{itemize}
\item one explicit toy model of debt accumulation across generations
\item one section translating the 8-register language into standard genetics language
\item one engagement section with Muller's ratchet, Hill-Robertson interference, and recombination advantage
\item one paragraph on what the register entries could physically mean
\end{itemize}

\textbf{Priority.} Very high, but draft after the toy model exists.

\subsection*{3. Toward a Zero-Junk Genome}
\textbf{Why this remains essential.} This is the platform paper. It should frame the entire DNA program and point to the sharper specialist papers.

\textbf{Current status.} Already drafted in \texttt{papers/tex/RS\_DNA\_Zero\_Junk\_Framework.tex}.

\textbf{What is stronger now than before.}
\begin{itemize}
\item full parity / degeneracy tables
\item explicit conservation hierarchy
\item helical phase stress shell
\item composite-cost interfaces
\item synthetic falsifier interfaces
\end{itemize}

\textbf{Second-pass improvements.}
\begin{itemize}
\item update the formalization summary to reflect the expanded genetics namespace
\item add a figure for the five-layer genome decomposition
\item add a codon-space figure
\item add an intron phase figure
\item move the most speculative ledger-balance language further into hypothesis framing
\end{itemize}

\textbf{Priority.} High, as the anchor preprint.

\subsection*{4. Z-Parity and the Hierarchy of Conservation in the Genetic Code}
\textbf{Why this moved up.} This paper is much more ready than it was on the first pass. The core theorem package exists and the story is clean.

\textbf{Lean base now in place.}
\begin{itemize}
\item \texttt{ZInvariantPipeline}
\item \texttt{AminoAcidZClassification}
\item \texttt{ZParityDegeneracyTable}
\item \texttt{WobbleZParity}
\item \texttt{ConservationHierarchy}
\end{itemize}

\textbf{Strongest honest claim.} Synonymous wobble substitutions can preserve amino-acid identity while breaking codon-level Z-parity, yielding a formal conservation hierarchy:
\[
\text{amino-acid identity} > \text{chemical similarity} > \text{Z-parity}.
\]

\textbf{Still needed before submission.}
\begin{itemize}
\item one compact visual table for all family classes
\item two or three examples beyond valine in the narrative section
\item one careful interpretation section on what parity-like structure could mean biologically
\end{itemize}

\textbf{Priority.} Write now.

\subsection*{5. Intron Positioning from Helical Phase Stress}
\textbf{Why this is important.} This is the most direct genomics test of the framework. If the predictor works, it is a real sequence-to-structure claim.

\textbf{Lean base now in place.}
\begin{itemize}
\item \texttt{IntronStructure}
\item \texttt{PhaseAlignmentLayer}
\item \texttt{HelicalPhaseModel}
\item \texttt{IntronPlacementOptimality}
\end{itemize}

\textbf{Strongest honest claim.} The theory now supports a dataset-independent notion of phase stress and candidate spacer insertion sites. What is missing is the empirical validation layer.

\textbf{Still needed before submission.}
\begin{itemize}
\item predictor implementation
\item curated benchmark genes
\item null-model comparisons
\item sensitivity analysis across transcript sets and annotation releases
\end{itemize}

\textbf{Priority.} One more strong validation pass.

\subsection*{6. The Molecular Gearbox: Phi-Ladder Timing in Biology}
\textbf{Why this matters.} This is the biophysics bridge paper linking the DNA program to water structure, molecular timing, and possibly circadian-scale dynamics.

\textbf{Lean base now in place.}
\begin{itemize}
\item \texttt{BioClocking}
\item \texttt{HydrationGearbox}
\item \texttt{GearboxRungs}
\item \texttt{GearboxSelectivity}
\item \texttt{Water/Constants}
\end{itemize}

\textbf{Strongest honest claim.} The current formal surface supports a phi-selective timing architecture and a rung-witness catalog. The mechanism is still partly interpretive and needs a tighter literature-facing section.

\textbf{Still needed before submission.}
\begin{itemize}
\item clean separation between theorem and mechanism
\item spectroscopy literature review
\item tighter physical argument for phi-selective filtering
\item more careful mapping from rungs to known biological timescales
\end{itemize}

\textbf{Priority.} One more mechanism-tightening pass.

\subsection*{7. Synthetic Biology Falsifiers for the RS DNA Framework}
\textbf{Why this should not wait too long.} This paper is strategically important because it converts the framework into an experimental program. The Lean layer now has explicit assay predicates and RS-violating construct interfaces.

\textbf{Lean base now in place.}
\begin{itemize}
\item \texttt{SyntheticFalsifierAssays}
\item \texttt{DegeneracySearchSpace}
\item \texttt{CodonAssignmentGlobalOptimality}
\item \texttt{IntronPlacementOptimality}
\item \texttt{CompositeJCost}
\end{itemize}

\textbf{Strongest honest claim.} We can now write a design paper describing construct classes and predicted assay failures, even before wet-lab collaboration exists.

\textbf{Still needed before submission.}
\begin{itemize}
\item concrete construct classes
\item assay definitions
\item expected qualitative effect directions
\item collaborator-facing experimental roadmap
\end{itemize}

\textbf{Priority.} Strategically high as a design / roadmap paper.

\subsection*{8. Deriving the Exact Codon-to-Amino-Acid Assignment}
\textbf{Why this is flagship.} If solved, this is the crown-jewel derivation: not just why there are 64 codons and 20 amino acids, but why the specific map is what it is.

\textbf{Lean base now in place.}
\begin{itemize}
\item \texttt{CodonAssignmentObjective}
\item \texttt{CodonAssignmentOptimality}
\item \texttt{WTokenAssignmentBridge}
\item \texttt{CompositeJCost}
\item \texttt{OptimizationCertificate}
\end{itemize}

\textbf{Why it is still not ready.} The theorem shell exists, but the composite cost is not yet calibrated well enough to support a strong biological claim.

\textbf{Still needed before submission.}
\begin{itemize}
\item a defensible composite metric
\item local optimality certificates under that metric
\item a clean bridge from WToken candidate classes to amino-acid classes
\item a careful statement of uniqueness versus near-uniqueness
\end{itemize}

\textbf{Priority.} Strategic flagship, not yet a full-paper claim.

\subsection*{9. Global J-Cost Optimality of the Genetic Code}
\textbf{Why this is still the hardest paper.} This is the strongest possible optimality paper and the least forgiving one. Reviewer scrutiny will be intense because the claim is global and the search space is enormous.

\textbf{Lean base now in place.}
\begin{itemize}
\item \texttt{DegeneracySearchSpace}
\item \texttt{CodonAssignmentGlobalOptimality}
\item \texttt{OptimizationCertificate}
\item \texttt{CompositeJCost}
\end{itemize}

\textbf{Still needed before submission.}
\begin{itemize}
\item combinatorial reductions that materially shrink the search space
\item certified exhaustive or bounded search over the reduced space
\item clear distinction between pattern-level and assignment-level optimality
\item a supplement explaining exactly what the certificates cover
\end{itemize}

\textbf{Priority.} Long-horizon flagship.

\section*{Recommended Writing Order}

\subsection*{Best tactical sequence for actual papers}
\begin{enumerate}
\item \textbf{The Standard Genetic Code Is Not Locally Optimal Under Simple Chemical Distance}
\item \textbf{Z-Parity and the Hierarchy of Conservation in the Genetic Code}
\item \textbf{Toward a Zero-Junk Genome} as revised anchor preprint
\item \textbf{Sex as Phase-Debt Cancellation}
\item \textbf{Intron Positioning from Helical Phase Stress}
\item \textbf{The Molecular Gearbox}
\item \textbf{Synthetic Biology Falsifiers} as a design paper
\end{enumerate}

\subsection*{Best flagship sequence}
\begin{enumerate}
\item \textbf{Deriving the Exact Codon-to-Amino-Acid Assignment}
\item \textbf{Global J-Cost Optimality of the Genetic Code}
\end{enumerate}

\textbf{Why this order now.}
\begin{itemize}
\item Paper 1 gives the sharpest clean result.
\item Paper 4 is small, legible, and helps build credibility fast.
\item Paper 3 anchors the full DNA program.
\item Paper 2 gives the strongest conceptual biology derivation once the toy model is added.
\item Papers 5 and 6 need one more validation / mechanism pass.
\item Papers 8 and 9 should only be written when the evidence package is genuinely ready.
\end{itemize}

\section*{Shared Assets To Build Once And Reuse Everywhere}
\begin{itemize}
\item a theorem ledger mapping claim, module, theorem name, and status
\item a common codon-space figure set
\item a common intron / helix figure set
\item deterministic optimizer result bundles
\item a standard literature comparison appendix
\item a shared glossary mapping RS language to standard biology language
\end{itemize}

\section*{Immediate Next Actions}
\textbf{Paper 1 package}
\begin{itemize}
\item freeze optimizer outputs
\item regenerate swap tables and baselines
\item draft methods and results first
\end{itemize}

\textbf{Paper 4 package}
\begin{itemize}
\item turn the new parity modules into one compact figure and one table
\item draft this as the fastest small paper
\end{itemize}

\textbf{Paper 3 package}
\begin{itemize}
\item revise \texttt{papers/tex/RS\_DNA\_Zero\_Junk\_Framework.tex}
\item add figures
\item update the formalization section to reflect the expanded Lean surface
\end{itemize}

\textbf{Paper 2 package}
\begin{itemize}
\item build one toy simulation of debt accumulation under asexual inheritance
\item write the translation section into mainstream recombination language
\end{itemize}

\textbf{Paper 5 package}
\begin{itemize}
\item implement a first intron predictor
\item benchmark against at least one null model
\end{itemize}

\textbf{Paper 6 package}
\begin{itemize}
\item separate proved statements from mechanistic conjectures more aggressively
\item tighten the spectroscopy and water-cluster section
\end{itemize}

\section*{Bottom-Line Recommendation}
If the goal is best immediate submission, do:
\begin{enumerate}
\item \textbf{Code not chemically optimal}
\item \textbf{Z-parity hierarchy}
\item \textbf{Zero-junk genome} as revised anchor
\end{enumerate}

If the goal is best conceptual impact, do:
\begin{enumerate}
\item \textbf{Code not chemically optimal}
\item \textbf{Sex as phase-debt cancellation}
\item \textbf{Zero-junk genome}
\end{enumerate}

If the goal is ultimate flagship value, keep working toward:
\begin{enumerate}
\item \textbf{Exact codon assignment}
\item \textbf{Global J-cost optimality}
\item \textbf{Synthetic falsifiers}
\end{enumerate}

The key second-pass conclusion is simple: the DNA paper program is now strong enough that the next bottleneck is not Lean formalization. It is paper architecture, deterministic computation, validation datasets, and reviewer-facing framing. The smaller and sharper papers should start now while the flagship derivation papers continue maturing in parallel.

\end{document}
